\documentclass[11pt]{article}
\usepackage[margin=1in]{geometry}
\usepackage{amsmath,amssymb}
\title{Exact Bregman splitting without an attained solution}
\author{ziangni-sys}
\date{4 October 2026}
\begin{document}
\maketitle
\noindent\textbf{Conjecture 00000008836.} The assertion that splitting under Bregman distances always converges is false without a solvability hypothesis. Both source versions state unconditional convergence. The following example is an exact Bregman proximal method, hence the zero-error special case of a relative-error method.

Work on the real Hilbert space $\mathbb R$, with the smooth mirror function $h(x)=x^2/2$ and the proper closed convex objective $f(x)=-x$. The second split term can be the zero function. The actual mirror gradient is $\nabla h(x)=x$, and its Bregman distance is
\[
 D_h(y,x)=h(y)-h(x)-\langle\nabla h(x),y-x\rangle
 =\frac{(y-x)^2}{2}.
\]
Thus $D_h\geq0$, with equality exactly when $y=x$. The mirror has strong-convexity modulus one, since its defining first-order lower bound holds with equality.

The genuine unit-penalty proximal subproblem is
\[
 \min_y Q_x(y),\qquad Q_x(y)=f(y)+D_h(y,x)
 =-y+\frac{(y-x)^2}{2}.
\]
For every current state $x$ and every competitor $z$,
\[
 Q_x(z)-Q_x(x+1)=\frac{(z-(x+1))^2}{2}.
\]
Consequently the unique global minimizer is $y=x+1$. All subproblems are well-defined and uniquely solvable, despite the lack of a minimizer of the original objective.

The actual convex subgradient relation is $\partial f(y)=\{-1\}$ at every $y$: its definition is the inequality $v(z-y)\leq f(z)-f(y)$ for every $z$. Choosing $z=y\pm1$ forces $v=-1$, which conversely satisfies every inequality. This graph is maximal monotone. Indeed any monotone extension containing a point $(a,b)$ must be monotone against $(a-1,-1)$ and $(a+1,-1)$; these give $b\geq-1$ and $b\leq-1$.

At $y=x+1$, choose $v=-1$ and enlargement error $e=0$. Then the exact Bregman inclusion holds:
\[
 v+\nabla h(y)-\nabla h(x)=0,\qquad v\in\partial f(y).
\]
In this Euclidean mirror, the usual relative-error certificate is also satisfied with zero residual, for every admissible $0\leq\sigma<1$:
\[
 v\in\partial_e f(y),\quad e\geq0,\quad
 |v+y-x|^2+2e\leq\sigma^2|y-x|^2.
\]
Here $\partial_e f(y)$ is defined by $v(z-y)-e\leq f(z)-f(y)$ for all $z$. Its left-hand residual is zero. The associated exact proximal update is precisely $x_{k+1}=x_k+1$.

For every starting point, induction gives $x_k=x_0+k$, which has no finite limit. If it converged, its consecutive difference would converge to zero, whereas $x_{k+1}-x_k=1$. Weak convergence also fails, as the identity is a continuous linear functional on this real Hilbert space.

\noindent\textbf{Scope.} The objective has no minimizer: $f(x+1)<f(x)$ for every $x$. Its subgradient has no zero. This is the missing solvability hypothesis in the unconditional claim. We do not refute Bregman convergence theorems that assume a solution exists, and do not separately assess the stated summability characterization.

\noindent\textbf{Formal verification.} Lean proves the actual proper closed convex functions, mirror gradient and Bregman distance, full subgradient relation and maximality, the complete proximal argmin relation, the zero-relative-error inclusion, all iterates, and failure of norm and weak convergence for every starting point.
\end{document}
